\section{Number of Longest Increasing Subsequences}
\label{sec:dp-number-of-lis}

\problemstatement{Given $\textit{arr}$, find the \textbf{number of
distinct} longest increasing subsequences (not just their shared
maximum length).}

\begin{patternbox}
This closing problem of the chapter pairs Section~\ref{sec:dp-lis}'s
ending-here \textit{dp} table with a second, parallel table
$\textit{count}[i]$: \textbf{how many} subsequences of the maximum
length achieve $\textit{dp}[i]$ -- updated using the same
``exists / count'' duality first introduced in the DP on Subsequences
chapter (there distinguishing whether a target sum is
\emph{achievable} versus \emph{how many ways} achieve it).
\end{patternbox}

\subsection*{Deriving the count update rule}

While filling $\textit{dp}[i]$ in the usual way, also track
$\textit{count}[i]$, initialized to $1$ (every index is trivially one
subsequence of length $1$, by itself). For each $j<i$ with
$\textit{arr}[j]<\textit{arr}[i]$:
\begin{itemize}
  \item If $\textit{dp}[j]+1 > \textit{dp}[i]$: a \textbf{strictly
        longer} chain through $j$ has just been discovered. Update
        $\textit{dp}[i] \gets \textit{dp}[j]+1$ and \textbf{reset}
        $\textit{count}[i] \gets \textit{count}[j]$ (every previous,
        shorter-chain count at $i$ is now obsolete -- it counted chains
        of a length that's no longer the best).
  \item If $\textit{dp}[j]+1 = \textit{dp}[i]$: this is \textbf{another
        equally good} way to reach the current best length at $i$.
        \textbf{Accumulate}: $\textit{count}[i] \mathrel{+}=
        \textit{count}[j]$.
\end{itemize}
After filling both tables, let $\textit{maxLen} = \max_i \textit{dp}[i]$.
The answer is $\sum_{i:\ \textit{dp}[i]=\textit{maxLen}}
\textit{count}[i]$ (every index whose best chain achieves the global
maximum length contributes its count).

\subsection*{C++ implementation}

\begin{lstlisting}
#include <bits/stdc++.h>
using namespace std;

int numberOfLIS(vector<int>& arr) {
    int n = arr.size();
    vector<int> dp(n, 1), count(n, 1);

    for (int i = 0; i < n; i++) {
        for (int j = 0; j < i; j++) {
            if (arr[j] < arr[i]) {
                if (dp[j] + 1 > dp[i]) {
                    dp[i] = dp[j] + 1;
                    count[i] = count[j];
                } else if (dp[j] + 1 == dp[i]) {
                    count[i] += count[j];
                }
            }
        }
    }

    int maxLen = *max_element(dp.begin(), dp.end());
    int total = 0;
    for (int i = 0; i < n; i++) {
        if (dp[i] == maxLen) total += count[i];
    }
    return total;
}

int main() {
    vector<int> arr = {1, 3, 5, 4, 7};
    cout << numberOfLIS(arr) << endl; // 2
    return 0;
}
\end{lstlisting}

\subsection*{Dry run}

\dryrun{Take $\textit{arr} = [1,3,5,4,7]$.}

\begin{center}
\begin{tabular}{c|c|c|c|c|c}
$i$ & 0 & 1 & 2 & 3 & 4 \\ \hline
$\textit{arr}[i]$ & 1 & 3 & 5 & 4 & 7 \\ \hline
$\textit{dp}[i]$ & 1 & 2 & 3 & 3 & 4 \\ \hline
$\textit{count}[i]$ & 1 & 1 & 1 & 1 & 2
\end{tabular}
\end{center}

At $i=3$ (value $4$): only $j=0,1$ qualify ($1<4$, $3<4$); the best is
through $j=1$ ($\textit{dp}[1]+1=3$), giving $\textit{dp}[3]=3$,
$\textit{count}[3]=\textit{count}[1]=1$. At $i=4$ (value $7$):
$j=2$ gives $\textit{dp}[2]+1=4$ (new max, reset
$\textit{count}[4]=\textit{count}[2]=1$); then $j=3$ \emph{also} gives
$\textit{dp}[3]+1=4$ (ties the current best, accumulate:
$\textit{count}[4] \mathrel{+}= \textit{count}[3] = 1+1=2$).
$\textit{maxLen}=4$, achieved only at $i=4$, contributing
$\textit{count}[4]=2$: the two LIS of length $4$ are $1,3,5,7$ and
$1,3,4,7$.

\begin{complexitybox}
\textbf{Time Complexity.} $O(n^2)$.
\textbf{Space Complexity.} $O(n)$ for \textit{dp} and \textit{count}.
\end{complexitybox}

\begin{takeawaybox}
This chapter closes exactly where the DP on Subsequences chapter's
``exists / count / optimize'' framing predicted it would: once a
length-optimizing ending-here table exists, a parallel counting table
updated by the same reset-on-strict-improvement,
accumulate-on-tie rule answers ``how many'' for free. The same two-table
pairing pattern is worth trying on any optimization DP whenever a
follow-up problem asks for a count instead of (or in addition to) the
optimum itself.
\end{takeawaybox}
